\section{Introduction}
\label{sec:introduction}

Let $Q_n$ be the graph with vertex set $\{0,1\}^n$ in which two vertices
are adjacent when they differ in exactly one coordinate. For a finite graph
$H$, let $R(H)$ be the least positive integer $M$ such that every red-blue
coloring of the edges of $K_M$ contains a monochromatic subgraph isomorphic
to $H$. The copy is not required to be induced.

\begin{theorem}\label{thm:main}
There is an absolute constant $C>0$ such that
\[
R(Q_n)\le C2^n
\qquad\text{for every integer }n\ge0.
\]
\end{theorem}

\subsection{The problem and its historical development}

In their 1975 paper, Burr and Erd\H{o}s called a family of graphs an
\emph{$L$-set} if one constant bounds the Ramsey number of every member
by that constant times its order~\cite[Section~1]{burr-erdos1975}.
The arboricity of a graph is the minimum number of forests into which its
edges can be partitioned. They conjectured that every family of uniformly
bounded arboricity is an $L$-set. A graph is \emph{$d$-degenerate} if every nonempty subgraph
has a vertex of degree at most $d$; uniformly bounded arboricity and
uniformly bounded degeneracy are equivalent boundedness conditions for
families~\cite[Section~3]{burr-erdos1975}. Burr and Erd\H{o}s raised the
cube question separately in Section~7, asking whether the cubes form an
$L$-set. Theorem~\ref{thm:main} resolves that question positively.

The cube question sits at a natural density boundary. The first-moment
lower bound used by Burr and Erd\H{o}s implies that a family with one linear
Ramsey constant must satisfy
$e(H)/|V(H)|=O(\log |V(H)|)$ as its graph orders grow, where $e(H)$ is
the number of edges~\cite[Lemma~4.3 and Section~7]{burr-erdos1975}.
For the cube,
$e(Q_n)/|V(Q_n)|=n/2=\tfrac12\log_2|V(Q_n)|$.
Thus cubes reach the logarithmic order of edges per vertex that this
necessary condition still permits. They test linear Ramsey growth beyond
families with bounded degeneracy, at the largest possible order of this
density parameter up to a constant factor.

For $n\ge1$, parity of the coordinate sum gives the connected cube two
bipartition classes of size $2^{n-1}$. The standard lower-bound coloring
with two blocks~\cite[Introduction]{conlon-fox-lee-sudakov2016} gives
$R(Q_n)\ge3\cdot2^{n-1}-1$. Indeed, for $n\ge2$, take blocks of sizes
$2^n-1$ and $2^{n-1}-1$, color edges inside each block red, and color
crossing edges blue. Each red component is too small for $Q_n$, while a
blue copy would need $2^{n-1}$ vertices in each block. The case $n=1$ is
immediate. Theorem~\ref{thm:main} therefore determines the order of
$R(Q_n)$. It gives no numerical value for $C$ and does not determine
whether $R(Q_n)/2^n$ converges or what its limiting value would be.

One line of progress established linear Ramsey growth when a sparsity
parameter is fixed. Chv\'atal, R\"odl, Szemer\'edi and Trotter proved in
1983 that graphs of any fixed maximum degree have Ramsey number linear in
their order, with a constant depending on the degree~\cite{chvatal-rodl-szemeredi-trotter1983}.
Lee later proved the
corresponding assertion for every fixed degeneracy, resolving the
bounded-degeneracy form of the Burr--Erd\H{o}s conjecture~\cite[Theorem~1.1]{lee2017}.
These results did not settle the cube question
because both the maximum degree and the degeneracy of $Q_n$ equal $n$:
their constants may grow with the dimension.

A second line sought quantitative estimates when the degree grows with
the graph. Early work on cubes includes Beck's 1983 paper~\cite{beck1983}.
Graham, R\"odl and Ruci\'nski proved that a bipartite graph $H$ of maximum
degree at most $\Delta\ge1$ satisfies
$R(H)<8(8\Delta)^\Delta |V(H)|$~\cite[Theorem~1]{graham-rodl-rucinski2001}.
For $Q_n$ with $n\ge1$, whose order is $m=2^n$, this gives $R(Q_n)<8(16n)^n$, or
$m^{O(\log\log m)}$. Shi then obtained a bound polynomial in $m$, using a
refinement of Kostochka and R\"odl's common-neighborhood lemma~\cite{kostochka-rodl2001,shi2001},
and later refined the estimate~\cite{shi2007}. This moved the cube problem from a growing exponent in
$m$ to a fixed exponent.

The next estimates approached and then improved on the quadratic scale.
For positive integers $n$, Fox and Sudakov obtained
$R(Q_n)\le n2^{2n+5}$~\cite[Corollary~1.3]{fox-sudakov2009}, and Conlon,
Fox and Sudakov obtained
$R(Q_n)\le2^{2n+6}$~\cite[Corollary~4.2]{conlon-fox-sudakov2016}.
Lee's Theorem~1.3 and the paragraph following it give
$R(Q_n)\le2^{2n}+n^2 2^n$ for all sufficiently large $n$~\cite{lee2017}.
Tikhomirov's exponential improvement is
$R(Q_n)\le2^{2n-cn+1}+2$ for an absolute $c>0$ and all sufficiently large
$n$~\cite[Corollary~1.2]{tikhomirov2024}.
In terms of $m=2^n$, these estimates pass from an extra logarithmic factor
above $m^2$, to a constant multiple of $m^2$, to $(1+o(1))m^2$, and
finally to an exponent below $2$. The present theorem reaches a constant
multiple of $m$.

\paragraph{Consequences for related graph families.}
The bound transfers to targets occupying a controlled fraction of an
ambient cube. If $H$ is a not necessarily induced subgraph of $Q_d$ and
$2^d\le A|V(H)|$ for a fixed $A\ge1$, monotonicity gives
$R(H)\le R(Q_d)\le C2^d\le CA|V(H)|$. The bound on the ambient order
controls the constant for the smaller target.

For example, let $d\ge0$ and $k\ge1$ be integers, and let $F$ be a
spanning subgraph of $Q_d$. Its \emph{Cartesian power} $F^{\square k}$
has vertex set $V(F)^k$; two tuples are adjacent when they agree in all
but one coordinate and the two entries there are adjacent in $F$.
Concatenating the $k$ blocks of cube coordinates makes this product a
spanning subgraph of $Q_{dk}$. Consequently,
\[
R(F^{\square k})\le R(Q_{dk})\le C2^{dk}=C|V(F)|^k,
\]
with the same absolute $C$ for all permitted $F,d,k$.

This includes Cartesian powers of paths on $2^d$ vertices for $d\ge1$,
and cycles on $2^d$ vertices for $d\ge2$. Starting with the order $0,1$
for $Q_1$, obtain an order of $Q_d$ by prefixing $0$ to the order of
$Q_{d-1}$ and then $1$ to its reverse. Consecutive vertices are adjacent,
and for $d\ge2$ the first and last vertices are also adjacent, giving the
cycle. Writing $P_a$ for the path on $a$ vertices, Mota, S\'ark\"ozy,
Schacht and Taraz proved
$R(P_a\square P_b)=(3/2+o(1))ab$ as $ab\to\infty$ for positive integers
$a,b$~\cite[Corollary~1.4]{mota-sarkozy-schacht-taraz2015}.
Their framework fixes the degree and does not give a constant uniform in
a growing number of Cartesian factors; the deduction above does so when
the factor spans a cube.

\subsection{Background on embedding and sampling}

The common-neighborhood method addresses a basic extension step in
bipartite embedding. Once one target class has been placed, each vertex in
the other class needs an unused common neighbor of its already placed
neighbors. Under their respective density and size assumptions,
Kostochka and R\"odl's deterministic lemma~\cite[Lemma~1]{kostochka-rodl2001}
and Sudakov's probabilistic argument~\cite[Section~2]{sudakov2003} find
large host sets in which every subset of a prescribed size has many common
neighbors; Sudakov explicitly drew on Gowers's earlier sampling idea~\cite[Section~4, Lemma~11]{gowers1998}.
Later proofs allow a controlled
collection of exceptional subsets and arrange the first embedding so that
target neighborhoods avoid them. Conlon uses a hypergraph packing
argument~\cite{conlon2009bipartite}; independently, Fox and Sudakov use an
auxiliary hypergraph embedding~\cite[Lemmas~2.1--2.2]{fox-sudakov2009}.
Conlon, Fox and Sudakov use local avoidance of exceptional subsets and
collisions~\cite[Section~4]{conlon-fox-sudakov2016}.
These refinements give bounds of the form $2^{O(\Delta)}m$ for bipartite
targets of order $m$ and maximum degree $\Delta$, but the exponential
dependence on $\Delta=n$ still costs a power of the cube order.
Lee's embedding records a penalty for small common neighborhoods,
controls its average over target neighborhoods, and embeds the most
constrained vertices first~\cite[Section~2]{lee2017}.

Tikhomirov~\cite[Sections~1, 5 and~6]{tikhomirov2024} identifies an obstruction
to mapping a whole cube bipartition uniformly into one prepared set.
His proof instead distinguishes dispersed common neighborhoods, a denser
subgraph, and a block structure suited to a facet-wise cube embedding.
These density and structural alternatives provide a useful methodological
comparison for the reductions below. His Theorem~1.1 is a density theorem
for one color: for all sufficiently large $n$, it embeds $Q_n$ in every
bipartite graph of density at least $1/2$ whose two parts each have at least
$2^{(2-c)n}$ vertices, for an absolute $c>0$. The present proof uses the
absence of a cube in both colors of a counterexample to derive discrepancy.

A distinct line studies $R(Q_n,K_s)$, the least order of a complete graph
whose every red-blue edge coloring contains a red $Q_n$ or a blue $K_s$.
Conlon, Fox, Lee and Sudakov proved $R(Q_n,K_s)\le c_s2^n$ for every fixed
$s\ge3$~\cite[Theorem~1.1]{conlon-fox-lee-sudakov2016}.
Fiz Pontiveros, Griffiths, Morris, Saxton and Skokan obtained the triangle
asymptotic $R(Q_n,K_3)=(1+o(1))2^{n+1}$~\cite[Theorem~1.1]{fiz-pontiveros-griffiths-morris-saxton-skokan2016},
and their theorem for a fixed clique gives
\[
R(Q_n,K_s)=(s-1)(2^n-1)+1
\]
for each fixed $s\ge3$ and all sufficiently large $n$~\cite[Theorem~1.1]{fiz-pontiveros-griffiths-morris-saxton-skokan2014}.
The cube tilings in the Conlon--Fox--Lee--Sudakov and triangle papers use
subcubes obtained by fixing initial coordinates, then control the edges
between adjacent pieces. The present proof uses that geometry with piece
sizes chosen from retained patch masses and with different control of
crossing edges. These conclusions for a fixed clique do not imply the diagonal
cube theorem, where the target in both colors grows with $n$.

The task in the present diagonal problem is to retain a constant
multiplicative loss when the host order is only a slowly diverging multiple
of $2^n$. We work with a hypothetical counterexample sequence having two
disjoint host sides of size $N$ each and $N/2^n\to\infty$, with no
prescribed rate of divergence.

\subsection{The four phases of the proof}
\label{sec:four-proof-phases}

A \emph{role} is a vertex of the cube being embedded, and a \emph{label}
is a vertex of the host graph. We map the even and odd roles to disjoint
host sides. The construction must preserve every cube edge in one color
and assign distinct labels to distinct roles. After an injective assignment
$f$ of the odd roles, the possible labels for an even role $v$ form the
common-neighbor set
\[
 L_v=\bigcap_{u\sim v}N_G(f(u)),
\]
where $G$ is the chosen color. Most embedding arguments below construct
probability laws $p_v$ supported on $L_v$ with
$\sum_v p_v(x)\le1$ for every host label $x$. Summing over any set of roles
then gives Hall's condition and hence distinct representatives. In the
last part of the proof, we instead construct random two-element subsets
of $L_v$ and bound directly the probability that any collection of these
sets violates Hall's condition.

To describe the reductions, a \emph{patch} will mean a pair of host sets,
one on each side, with specified sampling laws $\mu,\nu$; the required
density properties vary between the constructions. Write
\[
 d_G(\mu,\nu)=\Pr\{xy\in G\},\qquad x\sim\mu,\quad y\sim\nu
 \quad\text{independently}.
\]
The \emph{width} of $\mu$ is $\log(N\max_x\mu(x))$: a width bound $w$
allows no atom larger than $e^w/N$. In particular, the uniform law on
$M$ labels has width $\log(N/M)$. A discrepancy bound controls
$|d_G(\mu,\nu)-1/2|$ for every pair of laws within specified width budgets.
An orientation specifies which host side receives the first law; reversing
it interchanges the sides and their width budgets.

A patch property is \emph{available} if, for some fixed $\kappa>0$ and
all sufficiently large $n$, a patch satisfying it exists after every
removal of at most $\kappa N$ labels from each side. The patch may
change with the removed sets.

\begin{enumerate}[leftmargin=*]
\item \textbf{Reduce to a sequence and construct the probability tools.}
If the theorem fails, Section~\ref{sec:reduction-and-notation} gives
counterexamples with two disjoint host sides of size $N$, where
$n\to\infty$, $N/2^n\to\infty$, and $N\le n2^n$. The construction must
work without any prescribed rate for the middle limit.
Section~\ref{sec:general-conventions-and-probabilistic-tools} provides
mixtures of patch laws with small average label masses, locally determined
choices of random data, and ways to sample labels injectively while
controlling their joint probabilities.

\item \textbf{Use density gains or deduce discrepancy.}
By Lemma~\ref{lem:stabilization}, we may pass to a subsequence and
remove $o(N)$ labels from each side so that each tested patch property is
either absent or available. At prescribed power widths,
Section~\ref{sec:bias-versus-purity-at-power-widths} then shows that
available bias of at least a suitable negative power of $n$ forces an
available nearly monochromatic patch at doubled widths, unless a cube can
already be embedded. The broad-side and
small-grid constructions in
Sections~\ref{sec:a-constant-broad-high-neighborhood-case}--\ref{sec:a-small-power-key-grid}
then embed a cube under their respective concentration and high-degree
hypotheses. These implications, applied in both colors of a counterexample,
exclude bias throughout an initial range of widths: every permitted law
pair has density close to one half.
Sections~\ref{sec:asymmetric-pure-patches-under-2}--\ref{sec:jump-to-large-bias-only-at-linear-budget}
extend this range, allowing much greater concentration on one side than
on the other. They also exclude the full-dimensional cluster witnesses
of Section~\ref{sec:full-dimensional-cluster-patches}, in which many
labels have unusually large pairwise common neighborhoods.
Section~\ref{sec:the-two-deep-orientations-some-estimates} records the
resulting discrepancy bounds with a small linear width budget on either
side and a power width on the other. The cluster exclusion also bounds
the scales of the smaller cluster patches used next.

\item \textbf{Allocate the cube to patches and use the larger gains.}
Section~\ref{sec:patch-tiling-under-9} extracts disjoint patches and
retains a family with one color, one orientation, and one of the finitely
many constructions described there. It allocates the cube to these
patches by subcubes obtained by fixing initial coordinates, with sizes
proportional to the patch sizes. This is the geometry used in earlier
cube-versus-clique tilings~\cite{conlon-fox-lee-sudakov2016,fiz-pontiveros-griffiths-morris-saxton-skokan2016}.
The fixed coordinates identify the cube edges crossing between patches.

Smaller host sets and crossing edges reduce the available common
neighborhoods, so the construction needs a compensating gain. For
comparison, suppose that each of $n$ independent restrictions retains
each label with probability $1/2+\delta$. A set of size $M$ would then
leave an expected $M(1/2+\delta)^n$ labels. Relative to $N2^{-n}$, the factor is
$(M/N)(1+2\delta)^n$: the accumulated density surplus can compensate for
the smaller support. The proof obtains the needed gain either from edge
density or from common-neighbor density inside cluster patches, and
bounds the additional losses from crossing edges and dependent choices.
Section~\ref{sec:internal-solver-and-profiles-for-these-tilings}
constructs the laws used inside cluster patches. When the gain is large
enough, Section~\ref{sec:high-modes-of-the-tiling} assigns the odd roles
injectively and bounds the column sums of the even laws supported on
$L_v$. Hall's theorem then completes the embedding.

\item \textbf{Complete the cases with smaller gains.}
Sections~\ref{sec:low-modes-geometry-and-cell-states}
and~\ref{sec:initial-list-gates-and-resampling-estimates} first assign
the odd roles and preserve lists for the even roles. A fixed linear map
on cube words partitions the odd roles into syndrome classes; a chosen
$O(\log n)$ of these classes receive provisional labels. Their labels
are replaced later, one class at a time, from reserved host sets. Each
even role has at most one neighbor in each late class.
Section~\ref{sec:late-dynamics-and-its-transfer-problem} controls how
these replacements shorten its list and leaves two distinct candidates
adjacent to all its assigned odd neighbors. A joint bound on the
candidate pairs rules out Hall-deficient collections, giving distinct
even labels.
\end{enumerate}

Appendix~\ref{app:regime-table} collects the quantitative branches,
including their width, density, and loss budgets. It is a reference for
the parameter choices after the corresponding constructions have been
introduced.

\paragraph{The probability comparisons.}
A recurring construction samples a hidden $k$-tuple on the first host
side and requires neighboring odd labels to be adjacent to every entry.
Given these odd labels, the tuple's conditional distribution is supported
on the $k$th Cartesian power of their common neighborhood. A bound on its largest atom
therefore gives a lower bound on the size of that neighborhood. The tuple
can also influence which observations are selected. Its likelihood must
include that selection event, with the selection rule recomputed when the
tuple is varied. Lemma~\ref{lem:gated-posterior} formulates this comparison;
Section~\ref{sec:bias-versus-purity-at-power-widths} gives its first
application.

Injectivity introduces another dependence. The calibrated sampler in
Lemma~\ref{lem:near-product-injection} preserves individual marginals
exactly and bounds probabilities of several specified outputs by their
product probabilities times a controlled factor. Exact marginals are
needed in later cluster constructions whose scales may remain fixed.
The clock sampler in Lemma~\ref{lem:clock-sampling} also avoids specified
local failures and provides a joint upper bound. Summing such bounds
against nonnegative functions permits the common-neighbor and load
estimates to be used after injective sampling. The later sequential
assignment requires further conditional comparisons, developed where the
successive observations are defined in
Section~\ref{sec:late-dynamics-and-its-transfer-problem}.


Figure~\ref{fig:proof-roadmap} records the dependencies between the
reductions. In particular, the cluster exclusion is used both before the
deep discrepancy conclusion and during the later patch extraction.
\begin{figure}[p]
\centering
\begin{tikzpicture}[
  x=1cm,y=1cm,
  box/.style={draw=black!65,rounded corners=2pt,fill=black!2,
    align=center,font=\small,inner sep=6pt,text width=6.0cm},
  wide/.style={box,text width=12.7cm},
  arrow/.style={-{Stealth[length=2mm]},line width=.65pt,draw=black!75}
]
\node[wide] (bad) at (0,0)
  {Hypothetical counterexamples (Section~\ref{sec:reduction-and-notation}):\\
   $n\to\infty$, $C_n=N/2^n\to\infty$, $N\le n2^n$};
\node[wide] (small) at (0,-1.8)
  {Weighted densities are close to $1/2$ at small widths\\
   after excluding configurations that yield a cube\\
   (Sections~\ref{sec:bias-versus-purity-at-power-widths}--\ref{sec:a-small-power-key-grid})};
\node[box] (intermediate) at (-3.4,-4.0)
  {Density bounds extend to unequal\\
   widths on the two host sides\\
   (Sections~\ref{sec:asymmetric-pure-patches-under-2}--\ref{sec:intermediate-powers-of-bias})};
\node[box] (clusters) at (3.4,-4.0)
  {Full-dimensional cluster patches\\
   are excluded
   (Section~\ref{sec:full-dimensional-cluster-patches})};
\node[wide] (deep) at (0,-6.3)
  {Discrepancy holds at a small linear width on either side\\
   and a power width on the other\\
   (Sections~\ref{sec:jump-to-large-bias-only-at-linear-budget}--\ref{sec:the-two-deep-orientations-some-estimates})};
\node[wide] (tiling) at (0,-8.3)
  {Disjoint host patches of one color and orientation\\
   receive subcubes specified by initial coordinates\\
   (Sections~\ref{sec:patch-tiling-under-9}--\ref{sec:internal-solver-and-profiles-for-these-tilings})};
\node[box] (high) at (-3.4,-10.8)
  {Larger gains (Section~\ref{sec:high-modes-of-the-tiling})\\
   Odd labels are assigned injectively\\
   Even column sums give Hall};
\node[box] (low) at (3.4,-10.8)
  {Smaller gains (Sections~\ref{sec:low-modes-geometry-and-cell-states}--\ref{sec:late-dynamics-and-its-transfer-problem})\\
   Preserve lists through late assignments\\
   Joint pair bounds give Hall};
\node[wide] (cube) at (0,-13.3)
  {A monochromatic copy of $Q_n$: contradiction\\
   Hence $\sup_{n\ge0}R(Q_n)/2^n<\infty$};
\draw[arrow] (bad) -- (small);
\draw[arrow] (small.south) -- (0,-2.75) -| (intermediate.north);
\draw[arrow] (small.south) -- (0,-2.75) -| (clusters.north);
\draw[arrow] (intermediate.south) |- (0,-5.45) -- (deep.north);
\draw[arrow] (clusters.south) |- (0,-5.45) -- (deep.north);
\draw[arrow] (deep) -- (tiling);
\draw[arrow] (clusters.east) -- ++(.42,0) |- (tiling.east);
\draw[arrow] (tiling.south) -- (0,-9.3) -| (high.north);
\draw[arrow] (tiling.south) -- (0,-9.3) -| (low.north);
\draw[arrow] (high.south) |- (0,-12.05) -- (cube.north);
\draw[arrow] (low.south) |- (0,-12.05) -- (cube.north);
\end{tikzpicture}
\caption{Logical dependencies of the proof. Each reduction is carried out
along the hypothetical counterexample sequence, passing to subsequences
and removing a vanishing fraction of labels where needed. The separate
arrow from cluster exclusion to patch extraction records its second use:
it bounds the measured scales of the remaining cluster patches.
One common color and orientation are fixed before cube allocation.
The final estimates apply to every sequence $C_n\to\infty$, with no
prescribed rate of divergence.}
\label{fig:proof-roadmap}
\end{figure}


\Needspace{24\baselineskip}
\subsection{Terminology}

All constants and exponents are fixed before $n\to\infty$. A fixed
threshold for a small patch scale need not grow with $n$.

\begin{center}
\begin{tabular}{@{}p{.22\linewidth}p{.71\linewidth}@{}}
\toprule
Term or symbol & Meaning \\
\midrule
Position or role & A vertex of the cube being embedded. \\
Label & A vertex of the host graph. \\
$X,Y,N$ & Disjoint host sides, each initially of size $N$; later width
normalizations continue to use this original $N$. \\
$C_n$ & The ratio $N/2^n$, tending to infinity at an unspecified rate. \\
Width of $\mu$ & $\log(N\max_x\mu(x))$, with natural logarithms. \\
Patch and $M_i$ & In the final tiling, a pair of retained host sets of common
size $M_i$, used for a portion of the cube. Earlier patches are specified
by their laws. \\
Raw law & The specified law before the avoidance constraints of the stage
being discussed. \\
Gate & An event indicator included in a probability or likelihood to
restrict it to specified local conditions; it does not normalize the law. \\
Profile & A probability distribution on a finite set of permitted local
choices, chosen to balance expected label masses. \\
\bottomrule
\end{tabular}
\end{center}
